\subsection{Implikation und Äquivalenz}
\label{subsec:fundamentals_implication_and_equivalence}

Mit Implikationen und Äquivalenzen stellen wir logische Zusammenhänge zwischen
Aussagen her und können präzise beschreiben, wie die Gültigkeit einer Aussage
von einer anderen abhängt.

\begin{DefinitionBox}{Implikation}{fundamentals_implication}
  Die \emph{Implikation} \EnglishTerm{implication} zeigt an, welche Aussage aus einer anderen folgt:
  \[
    A\Rightarrow B
  \]
  bedeutet: „Wenn $A$ gilt, dann gilt auch $B$.“
\end{DefinitionBox}


\begin{RemarkBox}
  Eine Implikation gilt nicht notwendigerweise in beide Richtungen: Aus
  $A\Rightarrow B$ folgt im Allgemeinen nicht $B\Rightarrow A$.
\end{RemarkBox}

\begin{ExampleBox}{Implikation}{fundamentals_implication_example}
  Für eine ganze Zahl $x$ gilt die folgende Implikation:
  \[
    x=4
    \quad\Rightarrow\quad
    x \text{ ist gerade}.
  \]
  Dreht man den Implikationspfeil lediglich um, erhält man dagegen eine falsche
  Aussage:
    \begin{center}
      \begin{ClozeArea}[0.6\linewidth]{1.5cm}
      \[
           x \text{ ist gerade} \quad \nRightarrow \quad x=4.
      \]
      \end{ClozeArea}
    \end{center}
  Das Gegenbeispiel $x=2$ widerlegt die Aussage.
\end{ExampleBox}

\begin{DefinitionBox}{Äquivalenz}{fundamentals_equivalence}
  Die \emph{Äquivalenz} \EnglishTerm{equivalence}
  \[
    A\Leftrightarrow B
  \]
  bedeutet, dass sowohl $A\Rightarrow B$ als auch $B\Rightarrow A$ gilt. Das bedeutet, dass die Aussagen $A$ und $B$ äquivalent, also logisch gleichwertig sind. Sie sind entweder beide wahr oder beide falsch.
\end{DefinitionBox}

\begin{samepage}
Eine Wahrheitstabelle prüft alle möglichen Fälle; durch das Ausfüllen sehen wir,
unter welchen Wahrheitswerten eine Implikation oder eine Äquivalenz gilt.

\begin{center}
  \renewcommand{\arraystretch}{1.55}
  \setlength{\tabcolsep}{12pt}
  \begin{tabular}{cccc}
    \toprule
    $A$ & $B$ & $A\Rightarrow B$ & $A\Leftrightarrow B$\\
    \midrule
    T & T & \Cloze{1.8cm}{T} & \Cloze{1.8cm}{T}\\
    T & F & \Cloze{1.8cm}{F} & \Cloze{1.8cm}{F}\\
    F & T & \Cloze{1.8cm}{T} & \Cloze{1.8cm}{F}\\
    F & F & \Cloze{1.8cm}{T} & \Cloze{1.8cm}{T}\\
    \bottomrule
  \end{tabular}
\end{center}
\end{samepage}

Bei einer Implikation ist $A$ eine hinreichende Bedingung für $B$ und $B$ eine
notwendige Bedingung für $A$. Für eine Äquivalenz müssen diese Beziehungen in
beiden Richtungen gelten.
